% Part: lambda-calculus

\documentclass[../../include/open-logic-part]{subfiles}

\begin{document}

\olpart{lam}{లాంబ్డా కలనశాస్త్రం}

\begin{editorial}
ఈ భాగం లాంబ్డా కలనశాస్త్రం గురించి. పరిచయ అధ్యాయం జెరెమీ అవిగాడ్
గమనికల ఆధారంగా తయారైంది. అందులోని కొంత విషయాన్ని తరువాతి అధ్యాయాల్లోనూ వివరించాం.
అందువల్ల పరిచయంలో ఆ భాగం ఇప్పుడు అవసరం లేదు. వాక్యనిర్మాణం, చర్చ్--రోసర్ లక్షణం,
లాంబ్డా నిర్వచనీయతపై అధ్యాయాలను జెసెన్ చియన్ మిటాక్స్ వేసవి
ఇంటర్న్‌షిప్‌లో రూపొందించారు. వాటిని ఇంకా సమీక్షించి సవరించాలి.
\end{editorial}

\olimport[introduction]{introduction}

\olimport[syntax]{syntax}

\olimport[church-rosser]{church-rosser}

\olimport[lambda-definability]{lambda-definability}

\OLEndPartHook

\end{document}
